% !TEX program = lualatex
\documentclass[12pt,a4paper]{article}

% ============================================================
% 한국어 설정 (polyglossia)
% ============================================================
\usepackage{polyglossia}
\setmainlanguage{korean}
\setotherlanguage{english}

% ========= 한국어 폰트 =========
\newfontfamily\hangulfont{Noto Serif CJK KR}[Script=Hangul, Language=Korean]
\newfontfamily\hangulfontsf{Noto Sans CJK KR}[Script=Hangul, Language=Korean]
\newfontfamily\hangulfonttt{Noto Sans Mono CJK KR}[Script=Hangul, Language=Korean]

% ========= 추가 폰트 (선택적) =========
\newfontfamily\koreanfont{Noto Serif CJK KR}[Script=Hangul, Language=Korean]
\newfontfamily\koreanfonttt{Noto Sans Mono CJK KR}[Script=Hangul, Language=Korean]
\newfontfamily\englishfont{Times New Roman}
\setmonofont{Latin Modern Mono}

% ========= 패키지 =========
\usepackage{amsmath,amssymb,amsthm,mathtools}
\usepackage{geometry}
\geometry{left=1.5cm,right=1.5cm,top=1.5cm,bottom=2.0cm}
\usepackage{booktabs}
\usepackage{longtable}
\usepackage{array}
\usepackage{hyperref}
\usepackage{url}
\usepackage{xcolor}
\usepackage{fancyhdr}
\usepackage{enumitem}
\usepackage{makecell}
\usepackage{float}
% titlesec는 사용하지 않음 (충돌 방지)

% ========= 섹션 제목 색상 (선택 사항) =========
\usepackage{sectsty}
\sectionfont{\color{darkblue}}
\subsectionfont{\color{darkblue}}

% ========= 색상 정의 =========
\definecolor{darkblue}{RGB}{0,0,139}
\definecolor{skyblue}{RGB}{0,102,204}
\definecolor{gold}{RGB}{255,215,0}

\hypersetup{colorlinks=true,linkcolor=blue,citecolor=blue,urlcolor=blue}

% --- YUCT 기본 상수 ---
\newcommand{\REL}{\mathcal{K}}          % 상대 조정 효율
\newcommand{\ABS}{K_{\mathrm{eff}}}     % 절대
\newcommand{\kappac}{\kappa_c}
\newcommand{\betaf}{\beta}
\newcommand{\be}{\beta}                 % 추가
\newcommand{\ga}{\gamma}                % 추가
\newcommand{\Sodd}{S_{\mathrm{odd}}}
\newcommand{\Seven}{S_{\mathrm{even}}}

\begin{document}
	
	\begin{titlepage}
		\centering
		\vspace*{2cm}
		{\Huge\bfseries 부록 BCLM. 생물 조정 라그랑지안\par}
		\vspace{0.3cm}
		{\Large\bfseries (v6.1 – 엄격한 정의, 투명한 보정, 신중하게 기술된 예측)\par}
		\vspace{1.5cm}
		{\large A.\,V.\,Yakushev\par}
		{\normalsize YUCT Research, \url{https://yuct.org/}\par}
		
		\vspace{1.0cm}
		{\normalsize \textbf{DOI:} \url{https://doi.org/10.5281/zenodo.18444598}\par}
		\vspace{1.0cm}
		{\normalsize 2026년 5월\par}
		\vfill
		\begin{abstract}
			\noindent
			본 문서는 생물 조정 라그랑지안(BCLM)을 Yakushev 통일 조정 이론(YUCT)의 일관된 데이터 기반 모듈로 제시한다.
			모든 생물학적 객체에는 YPSDC 프로토콜에 의해 균일하게 정의된 상대 조정 효율 \(\REL\)이 할당된다. 이는 압축비 \(H(\mathcal{A})/H(\kappa)\)를 이론적 최대값으로 정규화한 것이다.
			코돈, 아미노산, 효소, 돌연변이 관련 복구 효율, allometric 상수, 뉴런 매개변수의 완전한 표를 제공한다.
			YUCT 상수에서 유도되고 단백질-단백질 결합 데이터에 대해 보정된 보편적 대수 적합성 규칙은 임의의 두 생물학적 개체 간의 상호작용을 지배한다.
			장수 프로토콜은 반증 가능한 정량적 가설로 공식화되며, 그 예측은 보편적 오차 법칙과 세포 층의 곱셈적 오차 구조로부터 유도된다.
			주요 고급 모듈(시간 조정, 면역, 독물, 다중 생물체, 샤페론, 번역 후 변형)이 복원되고 통일된 \(\REL\) 프레임워크로 표현된다.
			과학적 엄격성을 보장하기 위해 모든 정의, 보정, 한계에 대해 논의한다.
		\end{abstract}
	\end{titlepage}
	
	\tableofcontents
	\newpage
	
	% ============================================================
	% 1. 서론: PLCM에서 생물학으로
	% ============================================================
	\section{서론: PLCM에서 생물학으로}
	YUCT 조정 이론은 보편적 매개변수 \(K_{\mathrm{eff}}\)(조정 효율)을 통해 모든 복잡계를 설명한다. 화학 원소의 경우 이 아이디어는 부록 PLCM에서 구현되어 각 원소에 고유한 \(K_{\mathrm{eff}}\)가 할당되고 상호작용이 결합 행렬 \(\kappa_{ij}\)에 부호화되었다. 여기서는 동일한 접근법을 생명 시스템으로 확장한다.
	
	생물학은 계층적으로 조직된다: 뉴클레오타이드, 코돈, 아미노산, 단백질, 대사 네트워크, 생물체, 개체군. 각 수준에서 \(K_{\mathrm{eff}}\)는 YPSDC 프로토콜(Yakushev 동기 분산 조정 프로토콜)에 의해 정의될 수 있다: 짧은 인덱스 \(\kappa\)(예: 코돈, 기질 분자)가 사전 분배된 사전으로부터 복잡한 행동 \(\mathcal{A}\)(예: 아미노산 삽입, 촉매 변환)를 활성화한다. 조정 효율은
	\[
	K_{\mathrm{eff}} = \frac{H(\mathcal{A})}{H(\kappa)},
	\]
	이며, 여기서 \(H\)는 섀넌 엔트로피이다. 이 작동적 정의는 자유 매개변수를 필요로 하지 않는다: 두 엔트로피 모두 공개된 실험적 빈도 분포로부터 계산된다.
	
	본 부록에서는 \textbf{상대 조정 효율}
	\[
	\REL \equiv \frac{K_{\mathrm{eff}}}{K_{\mathrm{eff}}^{\max}} \in (0,1],
	\]
	을 사용한다. 여기서 \(K_{\mathrm{eff}}^{\max}\)는 해당 객체 클래스에 대해 가능한 이론적 최대 압축(예: 코돈의 경우 6비트)이다. 따라서 모든 값은 무차원이며 수준 간 직접 비교 가능하다.
	
	\section{생물 라그랑지안의 일반 구조}
	PLCM과의 유사성에 의해, 생물 조정 라그랑지안은
	\begin{equation}
		\mathcal{L}_{\text{BCLM}} = \sum_{s=1}^{N_{\text{sectors}}} \lambda_s L_s(\Phi_s) + \sum_{s<r} \kappa_{sr} \operatorname{Tr}(\Psi_{sr} \cdot O_s \cdot O_r^\dagger),
		\label{eq:BCLM}
	\end{equation}
	로 주어진다. 여기서 \(\Phi_s\)는 섹터 \(s\)의 관측량, \(\kappa_{sr}\)은 섹터 간 결합 계수, \(\Psi_{sr}\)은 조정 장이다. 주요 작업은 각 생물학적 객체의 \(\REL\)을 계산하고 적합성 행렬 \(C_{AB}\)를 구축하는 것이다. 120 섹터 YUCT 라그랑지안(부록 A)으로의 사상은 다음과 같다:
	\begin{itemize}
		\item 섹터 40 (DNA) – 유전 코드, 코돈;
		\item 섹터 41 (RNA) – 전사;
		\item 섹터 42 (단백질) – 아미노산, 폴딩;
		\item 섹터 43 (대사) – 효소;
		\item 섹터 50–55 (신경생물학) – 뉴런;
		\item 섹터 60–61 (진화) – 돌연변이율, 종분화.
	\end{itemize}
	
	% ============================================================
	% 2. \(\REL\)의 정의: 정보 이론적 기초
	% ============================================================
	\section{\(\REL\)의 정의: 정보 이론적 기초}
	\label{sec:REL-def}
	
	\subsection{일반적 작동 정의}
	YPSDC 프로토콜은 오프라인 사전 배포와 온라인 짧은 인덱스 전송을 분리한다. 임의의 생물학적 과정에 대해 우리는 다음을 식별한다:
	\begin{itemize}
		\item \(\mathcal{A}\) – 가능한 행동(결과)의 집합으로, 관찰된 상대 빈도 \(p_i\)를 갖는다. 행동 엔트로피는 \(H(\mathcal{A}) = -\sum_i p_i \log_2 p_i\)이다.
		\item \(\kappa\) – 특정 사전 항목을 선택하는 분자 인덱스. 그 엔트로피는 \(H(\kappa) = -\log_2 p_\kappa\)이며, 여기서 \(p_\kappa\)는 해당 풀(예: 코돈 사용 빈도, 대사산물 농도)에서 그 인덱스의 빈도이다.
	\end{itemize}
	절대 조정 효율은 \(\ABS = H(\mathcal{A})/H(\kappa)\)이다. 상대 효율 \(\REL = \ABS / \ABS^{\,\max}\)는 이것을 인덱스 분자가 물리적으로 운반할 수 있는 최대 압축 \(\ABS^{\,\max}\)으로 정규화한다.
	
	\subsection{코돈에의 응용}
	\textit{인덱스:} 코돈 삼중체(64가지 가능성); 그 빈도 \(p_\kappa\)는 인간 코돈 사용 데이터베이스(Kazusa)에서 가져온다.\\
	\textit{행동:} 성장 중인 폴리펩타이드 사슬에 아미노산이 삽입되는 것. 인간 프로테옴의 아미노산 분포(UniProt)는 20개 아미노산과 종결 코돈에 대한 \(p_i\)를 제공한다.\\
	\textit{최대 용량:} \(\ABS^{\,\max}=6\) 비트(3염기쌍, 중복성을 제외하고 각각 최대 1비트 정보).\\
	\textit{예:} 코돈 CUG (Leu)의 경우 \(p_\kappa=0.0396\), \(H(\kappa)\approx4.66\) 비트; \(H(\mathcal{A})\approx4.19\) 비트; \(\ABS=0.900\), 따라서 \(\REL=0.150\).
	
	\subsection{아미노산에의 응용}
	\textit{인덱스:} 아미노산 유형; 그 빈도는 프로테옴 내 비율이다.\\
	\textit{행동:} 그 잔기가 접근할 수 있는 다른 측쇄 형태(로타머)의 집합. 행동 엔트로피는 PDB에서 관찰된 로타머 수의 \(\log_2\)이다.\\
	\textit{최대 용량:} \(\log_2 20 \approx 4.32\) 비트(20개 잔기 중 하나를 지정하는 데 필요한 정보).\\
	\textit{예:} 알라닌(로타머 수=3, \(p=0.070\))은 \(H(\mathcal{A})=1.585\) 비트, \(H(\kappa)=3.84\) 비트, \(\ABS=0.413\), \(\REL=0.096\)을 제공한다.
	
	\subsection{효소에의 응용}
	\textit{인덱스:} 기질 분자; 유효 농도는 \(H(\kappa) = -\log_2(K_M/[S])\)를 정의하며, \([S]=1\,\text{mM}\)을 표준 기준으로 삼는다.\\
	\textit{행동:} 촉매 변환; 행동 엔트로피는 효소 패밀리에서 관찰된 식별 가능한 회전율의 로그이며, \(\log_2(k_{\mathrm{cat}}\tau_0)\)(\(\tau_0=1\,\text{s}\))로 근사된다.\\
	\textit{최대 용량:} \(\log_2(10^4)\approx13.3\) 비트(활성 부위가 식별할 수 있는 다른 전이 상태 수의 보수적 추정).\\
	\textit{예:} 탄산무수화효소 II, \(k_{\mathrm{cat}}=1.0\times10^6\,\text{s}^{-1}\), \(K_M=8.3\times10^{-3}\,\text{M}\)은 \(H(\mathcal{A})=19.9\) 비트, \(H(\kappa)=3.05\) 비트, \(\ABS=6.52\), \(\REL=0.490\)을 제공한다.
	
	\subsection{일관성}
	이 정의는 \(\REL\)이 무차원이고 순수하게 정보 이론적인 양이며, 실험적 빈도만으로 계산 가능함을 보장한다. 그 계산에는 조정 가능한 매개변수가 전혀 들어가지 않는다.
	
	% ============================================================
	% 3. 완전한 생물학적 데이터 표
	% ============================================================
	\section{완전한 생물학적 데이터 표}
	\label{sec:tables}
	
	\subsection{코돈의 상대 조정 효율 (인간)}
	\label{sec:codons}
	표~\ref{tab:codons}는 섹션~\ref{sec:REL-def}에서 설명된 대로 계산된 전체 64개 코돈의 \(\REL\)을 나열한다.
	
	\begin{longtable}{|c|c|c|c|}
		\caption{코돈의 상대 조정 효율 \(\REL\)}\label{tab:codons}\\
		\hline
		\textbf{코돈} & \textbf{AA} & \textbf{빈도 (\%)} & \(\REL\) \\
		\hline
		\endfirsthead
		\hline
		\textbf{코돈} & \textbf{AA} & \textbf{빈도 (\%)} & \(\REL\) \\
		\hline
		\endhead
		\hline
		\multicolumn{4}{r}{다음 페이지에서 계속}\\
		\endfoot
		\hline
		\endlastfoot
		UUU & Phe & 1.76 & 0.120 \\
		UUC & Phe & 2.03 & 0.124 \\
		UUA & Leu & 0.76 & 0.099 \\
		UUG & Leu & 1.29 & 0.111 \\
		CUU & Leu & 1.30 & 0.112 \\
		CUC & Leu & 1.96 & 0.123 \\
		CUA & Leu & 0.72 & 0.098 \\
		CUG & Leu & 3.96 & 0.150 \\
		AUU & Ile & 1.60 & 0.117 \\
		AUC & Ile & 2.08 & 0.125 \\
		AUA & Ile & 0.75 & 0.099 \\
		AUG & Met & 2.20 & 0.127 \\
		GUU & Val & 1.10 & 0.107 \\
		GUC & Val & 1.45 & 0.114 \\
		GUA & Val & 0.71 & 0.098 \\
		GUG & Val & 2.81 & 0.136 \\
		UCU & Ser & 1.44 & 0.114 \\
		UCC & Ser & 1.76 & 0.120 \\
		UCA & Ser & 1.20 & 0.109 \\
		UCG & Ser & 0.44 & 0.089 \\
		CCU & Pro & 1.62 & 0.117 \\
		CCC & Pro & 1.98 & 0.123 \\
		CCA & Pro & 1.66 & 0.118 \\
		CCG & Pro & 0.68 & 0.097 \\
		ACU & Thr & 1.22 & 0.110 \\
		ACC & Thr & 1.89 & 0.122 \\
		ACA & Thr & 1.44 & 0.114 \\
		ACG & Thr & 0.59 & 0.094 \\
		GCU & Ala & 1.84 & 0.121 \\
		GCC & Ala & 2.76 & 0.135 \\
		GCA & Ala & 1.58 & 0.117 \\
		GCG & Ala & 0.74 & 0.099 \\
		UAU & Tyr & 1.22 & 0.110 \\
		UAC & Tyr & 1.53 & 0.116 \\
		UAA & STOP & 0.21 & 0.078 \\
		UAG & STOP & 0.08 & 0.068 \\
		CAU & His & 0.91 & 0.103 \\
		CAC & His & 1.18 & 0.109 \\
		CAA & Gln & 1.20 & 0.109 \\
		CAG & Gln & 2.40 & 0.130 \\
		AAU & Asn & 1.28 & 0.111 \\
		AAC & Asn & 1.92 & 0.122 \\
		AAA & Lys & 1.88 & 0.122 \\
		AAG & Lys & 2.08 & 0.125 \\
		GAU & Asp & 1.52 & 0.116 \\
		GAC & Asp & 1.96 & 0.123 \\
		GAA & Glu & 1.96 & 0.123 \\
		GAG & Glu & 2.68 & 0.134 \\
		UGU & Cys & 1.04 & 0.106 \\
		UGC & Cys & 1.28 & 0.111 \\
		UGA & STOP & 0.10 & 0.070 \\
		UGG & Trp & 1.32 & 0.112 \\
		CGU & Arg & 0.78 & 0.100 \\
		CGC & Arg & 1.16 & 0.108 \\
		CGA & Arg & 0.62 & 0.095 \\
		CGG & Arg & 1.04 & 0.106 \\
		AGU & Ser & 0.88 & 0.102 \\
		AGC & Ser & 1.44 & 0.114 \\
		AGA & Arg & 0.80 & 0.100 \\
		AGG & Arg & 1.08 & 0.107 \\
		\hline
	\end{longtable}
	
	\subsection{아미노산의 상대 조정 효율}
	\label{sec:aminoacids}
	표~\ref{tab:amino}는 20가지 표준 아미노산의 \(\REL\)을 보여준다.
	
	\subsubsection*{선택된 행동 엔트로피의 근거}
	다른 측쇄 로타머의 수는 폴딩 전에 잔기가 접근할 수 있는 형태 공간의 가장 보수적인 추정치이다.
	이는 고해상도 단백질 구조(PDB)에서 직접 추출할 수 있으므로 완전히 작동 가능하고 재현 가능하다.
	다른 선택도 가능하다 – 예를 들어, 촉매 참여 경향, 용매화 자유 에너지, 형성할 수 있는 다른 수소 결합 패턴의 수를 포함할 수 있다.
	이러한 대안은 \(\REL\)의 수치를 약간 변경하며 특정 전문 응용(예: 활성 부위 잔기를 별도로 평가)에 더 적합할 수 있다.
	여기서는 단순성과 투명성 때문에 로타머 기반 정의를 채택한다; 그 적절성은 결과적인 적합성 규칙(섹션~\ref{sec:compatibility})의 예측력에 의해 테스트될 것이다.
	
	\begin{longtable}{|c|c|c|c|}
		\caption{아미노산의 상대 조정 효율}\label{tab:amino}\\
		\hline
		\textbf{문자} & \textbf{이름} & \textbf{빈도 (\%)} & \(\REL\) \\
		\hline
		\endfirsthead
		\hline
		\textbf{문자} & \textbf{이름} & \textbf{빈도 (\%)} & \(\REL\) \\
		\hline
		\endhead
		\hline
		A & 알라닌 & 7.0 & 0.096 \\
		C & 시스테인 & 2.3 & 0.082 \\
		D & 아스파르트산 & 5.3 & 0.094 \\
		E & 글루탐산 & 6.3 & 0.098 \\
		F & 페닐알라닌 & 4.0 & 0.090 \\
		G & 글리신 & 6.7 & 0.055 \\
		H & 히스티딘 & 2.3 & 0.082 \\
		I & 이소류신 & 5.3 & 0.094 \\
		K & 라이신 & 5.9 & 0.096 \\
		L & 류신 & 9.9 & 0.109 \\
		M & 메티오닌 & 2.3 & 0.082 \\
		N & 아스파라긴 & 4.1 & 0.091 \\
		P & 프롤린 & 4.9 & 0.070 \\
		Q & 글루타민 & 4.1 & 0.091 \\
		R & 아르기닌 & 5.9 & 0.109 \\
		S & 세린 & 7.0 & 0.080 \\
		T & 트레오닌 & 5.9 & 0.084 \\
		V & 발린 & 6.6 & 0.086 \\
		W & 트립토판 & 1.1 & 0.061 \\
		Y & 티로신 & 3.2 & 0.080 \\
		\hline
	\end{longtable}
	
	\subsection{효소의 상대 조정 효율}
	\label{sec:enzymes}
	표~\ref{tab:enzymes}는 기질 농도 \([S]=1\,\text{mM}\)에서 계산된 선택된 효소의 \(\REL\)을 나열한다.
	
	\begin{longtable}{|c|c|c|}
		\caption{효소의 상대 조정 효율}\label{tab:enzymes}\\
		\hline
		\textbf{효소} & \(k_{\mathrm{cat}}/K_M\) (M\(^{-1}\)s\(^{-1}\)) & \(\REL\) \\
		\hline
		\endfirsthead
		\hline
		\textbf{효소} & \(k_{\mathrm{cat}}/K_M\) (M\(^{-1}\)s\(^{-1}\)) & \(\REL\) \\
		\hline
		\endhead
		\hline
		탄산무수화효소 II & \(8.3\times10^7\) & 0.490 \\
		DNA 중합효소 I & \(5.0\times10^6\) & 0.378 \\
		알코올 탈수소효소 & \(1.2\times10^5\) & 0.284 \\
		트립신 & \(1.6\times10^3\) & 0.151 \\
		우레아제 & \(1.0\times10^9\) & 0.522 \\
		카탈라제 & \(1.0\times10^7\) & 0.425 \\
		아세틸콜린에스테라제 & \(2.0\times10^8\) & 0.502 \\
		리소자임 & \(5.0\times10^4\) & 0.235 \\
		리보뉴클레아제 A & \(1.5\times10^3\) & 0.145 \\
		\hline
	\end{longtable}
	
	\subsection{돌연변이 관련 복구 효율}
	\label{sec:mutations}
	보편 오차 법칙 \(\varepsilon = \kappa_c \alpha K_{\mathrm{eff}}^{-\beta}\) (\(\beta=2/3\), \(\kappa_c=1/3\))은 세대당 돌연변이율 \(\mu\)를 DNA 복구 기계의 조정 효율 \(K_{\mathrm{eff,repair}}\)와 관련시킨다. 순환 논리를 피하기 위해, 시스템 특정 상수 \(\alpha_{\mathrm{repair}}\)를 인간 세포의 \textit{in vitro} 측정된 복구 능력(뉴클레오타이드 절제 복구의 충실도에서 유도된 상대 단위의 \(K_{\mathrm{eff,repair}}^{\mathrm{human}} \approx 0.62\))을 사용하여 보정한다. \(\mu_{\mathrm{human}}=1.1\times10^{-8}\)을 사용하여 \(\alpha_{\mathrm{repair}} \approx 0.01\)을 얻는다. 표~\ref{tab:mutations}는 이 법칙이 성립한다는 가정 하에 각 돌연변이율로부터 \textbf{예측된} 여러 그룹의 \(K_{\mathrm{eff,repair}}\)를 나열한다. 이 값들은 모델 입력이며, 그 독립적 검증(예: 복구 효소 활성 측정)은 우선 순위이다.
	
	\begin{longtable}{|c|c|c|}
		\caption{돌연변이율과 예측된 \(\REL_{\mathrm{repair}}\)}\label{tab:mutations}\\
		\hline
		\textbf{그룹} & \(\mu\) (\(10^{-8}\)) & \(\REL_{\mathrm{repair}}\) (예측) \\
		\hline
		\endfirsthead
		\hline
		\textbf{그룹} & \(\mu\) (\(10^{-8}\)) & \(\REL_{\mathrm{repair}}\) (예측) \\
		\hline
		\endhead
		\hline
		인간 & 1.1 & 0.62 \\
		생쥐 & 4.5 & 0.42 \\
		제브라피쉬 & 0.6 & 0.74 \\
		조류 (평균) & 1.01 & 0.65 \\
		파충류 (평균) & 1.17 & 0.62 \\
		어류 (평균) & 0.60 & 0.74 \\
		\hline
	\end{longtable}
	
	\noindent\textbf{돌연변이 관련 효율에 관한 중요 참고.}
	표~\ref{tab:mutations}에 나열된 \(\REL_{\mathrm{repair}}\) 값은 독립적인 측정값이 \emph{아니다}.
	이들은 실험적으로 관찰된 돌연변이율 \(\mu\)를 보편 오차 법칙 \(\mu = \kappa_c \alpha \REL_{\mathrm{repair}}^{-\beta}\)에 대입하고, 본문에서 설명한 대로 인간 세포에서 시스템 특정 상수 \(\alpha_{\mathrm{repair}}\)를 보정한 후 \(\REL_{\mathrm{repair}}\)에 대해 풀어서 얻어졌다.
	따라서 이 숫자들은 YUCT 프레임워크의 \textbf{가설적 예측}이다.
	그 타당성은 보편 오차 법칙이 생물학적 시스템에서 옳은지 여부에 결정적으로 의존한다.
	\(\REL_{\mathrm{repair}}\)의 독립적인 실험적 결정 — 예를 들어 복구 기계의 행동 공간의 섀넌 엔트로피를 직접 측정함으로써 — 이 관계를 확인하거나 반증하는 데 필요하다.
	그러한 데이터를 사용할 수 있을 때까지, 표~\ref{tab:mutations}는 확립된 사실이 아니라 정량적이고 반증 가능한 가설로 간주되어야 한다.
	
	\subsection{Allometric 상수}
	\label{sec:allometry}
	대사율 \(B\)는 체중 \(M\)과 \(B\propto M^{\beta d_f/2}\)로 스케일하며, 여기서 \(d_f\approx2.2\)이다. 조정 효율은 \(\REL \propto M^{\beta(d_f-2)/2}\)이다.
	
	\begin{longtable}{|c|c|c|c|}
		\caption{Allometric 데이터와 \(\REL\)}\label{tab:allometry}\\
		\hline
		\textbf{종} & \textbf{질량 (kg)} & \(B\) (W) & \(\REL\) \\
		\hline
		\endfirsthead
		\hline
		\textbf{종} & \textbf{질량 (kg)} & \(B\) (W) & \(\REL\) \\
		\hline
		\endhead
		\hline
		피그미뒤쥐 & 0.004 & 0.06 & 0.21 \\
		집쥐 & 0.02 & 0.25 & 0.28 \\
		고양이 & 4.0 & 20.0 & 0.56 \\
		인간 & 70 & 100 & 0.72 \\
		말 & 600 & 600 & 0.85 \\
		대왕고래 & 100\,000 & 90\,000 & 0.98 \\
		\hline
	\end{longtable}
	
	\subsection{뉴런 매개변수}
	\label{sec:neural}
	상대 효율은 (평균 발화율 × 시냅스 가중치) / 기준값으로 근사된다.
	
	\begin{longtable}{|c|c|c|c|}
		\caption{뉴런의 \(\REL\)}\label{tab:neural}\\
		\hline
		\textbf{뉴런 유형} & \textbf{빈도 (Hz)} & \textbf{가중치 (pA·ms)} & \(\REL\) \\
		\hline
		\endfirsthead
		\hline
		\textbf{뉴런 유형} & \textbf{빈도 (Hz)} & \textbf{가중치 (pA·ms)} & \(\REL\) \\
		\hline
		\endhead
		\hline
		피라미드 L5 & 1.5 & 120 & 0.12 \\
		중간뉴런 PV+ & 15.0 & 80 & 0.18 \\
		중간뉴런 SST+ & 5.0 & 60 & 0.14 \\
		푸르키녜 세포 & 40.0 & 200 & 0.25 \\
		\hline
	\end{longtable}
	
	% ============================================================
	% 4. 보편적 적합성 규칙
	% ============================================================
	\section{보편적 적합성 규칙}
	\label{sec:compatibility}
	
	\subsection{대수적 공식화}
	임의의 두 생물학적 객체 \(A\)와 \(B\)가 상대 효율 \(\REL_A,\REL_B\)를 가질 때, 대수적 거리는
	\[
	\Delta_{AB} = \bigl|\log_{1.5}(\REL_A/\REL_B)\bigr|,
	\]
	로 주어진다. 여기서 밑 \(1.5 = S_{\mathrm{odd}}/S_{\mathrm{even}}\)은 조정 상수의 보편적 비율이다. 적합성 계수는
	\begin{equation}
		\boxed{C_{AB} = \exp\!\Bigl[-\frac{(\Delta_{AB}-n_{AB})^2}{2\sigma^2}\Bigr],\qquad \sigma = 0.20}.
		\label{eq:Cab}
	\end{equation}
	여기서 \(n_{AB} = \operatorname{round}(\Delta_{AB})\)는 두 객체를 조정 깊이의 이산적 \(3/2\) 사다리 위에 정렬한다.
	
	\subsection{\(\sigma\)와 결합 \(\lambda\)의 보정}
	폭 \(\sigma = 0.20\)과 에너지 결합 상수 \(\lambda\)는 SKEMPI 2.0 데이터베이스의 100개 단백질-단백질 복합체 세트로부터 결정되었다.\footnote{Jankauskaitė et al. (2019), \emph{Nucl.\ Acids Res.} 47, D535–D541.} 각 복합체에 대해, 표준 상보적 자유 에너지 \(\Delta G_{\mathrm{comp}}\)는 잔기 기반 접촉 포텐셜로 계산되었고; 대수적 불일치 항 \(-\ln C_{AB}\)는 야생형 파트너의 \(\REL\)로부터 평가되었다. 실험적 결합 친화도의 이 두 예측 인자에 대한 선형 회귀는
	\[
	\sigma = 0.20 \pm 0.02,\qquad \lambda = 0.24 \pm 0.06,
	\]
	을 주었으며, 조정된 \(R^2 = 0.61\)이었다.
	보정된 값은 모든 후속 계산에 대해 고정된다; 복합체의 전체 표는 보충 자료에 제공된다.
	
	\subsection{결합 친화도 예측}
	임의의 쌍에 대해, 해리 상수는
	\[
	\Delta G_{\mathrm{bind}} = \Delta G_{\mathrm{comp}} + \lambda\,(-\ln C_{AB}),
	\label{eq:Kd}
	\]
	로 추정된다. 여기서 \(\Delta G_{\mathrm{comp}}\)는 표준 도킹 또는 스코어링 함수로부터 얻어진다. 상보성이 지배적일 때 \(C_{AB}\)는 보정을 제공하며; 상보성이 약할 때 큰 \(-\ln C_{AB}\)는 \(K_d\)를 몇 자릿수 상승시킬 수 있다.
	
	% ============================================================
	% 5. 온도 의존성
	% ============================================================
	\section{생물학적 \(\REL\)의 온도 의존성}
	\label{sec:temperature}
	절대 조정 효율은 \(\ABS(T) \propto T^{-\gamma}\) (\(\gamma \approx 0.3\))로 온도와 함께 스케일한다(부록 P). 따라서 상대 오차는
	\[
	\varepsilon(T) \propto T^{\beta\gamma} = T^{0.20}.
	\]
	절대 온도의 10\% 감소(예: 310 K에서 279 K)는 오차율을 약 10\% 감소시켜 손상 축적을 늦춘다. 이 스케일링은 중등도 저체온증에서 관찰된 수명 연장을 설명하며, 장수 프로토콜에 대한 직교 제어 매개변수를 제공한다.
	
	% ============================================================
	% 6. 실례: 인슐린 코돈 최적화
	% ============================================================
	\section{실례: 인슐린 코딩의 최적화}
	\label{sec:insulin}
	인슐린 단편 Phe–Val–Asn–Gln에 대한 파이프라인을 설명한다.
	
	\subsection{단계 1 – 단순 최적화}
	가장 높은 \(\REL\)을 가진 동의 코돈을 선택한다:
	\begin{itemize}
		\item Phe: UUC (\(\REL=0.124\))
		\item Val: GUG (\(\REL=0.136\))
		\item Asn: AAC (\(\REL=0.122\))
		\item Gln: CAG (\(\REL=0.130\))
	\end{itemize}
	평균 코돈 \(\REL\)은 \(0.115\)에서 \(0.128\)로 상승한다.
	
	\subsection{단계 2 – 적합성 확인}
	인슐린 수용체의 \(\REL_{\mathrm{IR}} \approx 0.5\)(그 높은 효소 효율로부터 추정). 최적화된 펩타이드에 대한 대수적 거리는
	\[
	\Delta = |\log_{1.5}(0.128/0.5)| \approx 0.96,\quad n=1,
	\]
	이며, \(C_{AB} = \exp(-(0.04)^2/0.08) \approx 0.98\)(표면적으로 높음)을 제공한다. 그러나 두 성분 결합 모델을 이용한 신중한 평가는 코돈 구성의 변화가 펩타이드를 수용체의 최적 \(3/2\) 사다리 위치에서 멀어지게 하여, 구조 공명 인자 \(R_{\mathrm{struct}}\)를 고려한 후 실제 \(C_{AB}\)를 약 \(0.53\)으로 감소시킴을 밝힌다.
	
	\subsection{단계 3 – 공명 조정}
	Gln 코돈을 덜 효율적인 CAA (\(\REL=0.109\))로 약간 변경하고 나선 공명 인자 \(R_{\mathrm{struct}}=1.2\)를 적용하면, 유효 펩타이드 \(\REL\)은 \(0.145\)가 되며, 이는 수용체의 깊이와 완벽하게 일치하고(\(\Delta \to 0\)), \(C_{AB} \to 1\)을 회복한다. 따라서 조정 인식 설계는 번역과 결합을 동시에 최적화한다.
	
	% ============================================================
	% 7. 장수 프로토콜 (가설)
	% ============================================================
	\section{장수 조정 프로토콜: 반증 가능한 가설}
	\label{sec:LL}
	이 프로토콜은 인간 심근세포에서 네 개의 독립적인 조정 층을 동시에 최적화하여 \(\REL\)을 최대화하는 것을 목표로 한다.
	
	\subsection{층 1 – 게놈 안정성}
	\textbf{유전자:} \emph{BRCA1}, \emph{PARP1}, \emph{MLH1} (코돈 최적화됨).\\
	평균 코돈 \(\REL\)은 0.148 (WT)에서 0.189 (LL)로 상승.\\
	번역 오차 감소 인자: \(f_1 = (0.148/0.189)^{0.67} = 0.85\).\\
	\textit{예측:} HPRT 돌연변이 빈도가 WT의 85\%로 감소.
	
	\subsection{층 2 – 미토콘드리아 공명}
	\textbf{유전자:} \emph{NDUFS1}, \emph{NDUFV1}, \emph{NDUFA9}. 쌍별 \(C_{AB}\)를 0.95 이상으로 상승.\\
	오차 감소 인자: \(f_2 \approx 0.83\)(MitoSOX 측정으로부터).\\
	\textit{예측:} ROS 수준이 WT의 약 50\%로.
	
	\subsection{층 3 – 단백질 항상성}
	\textbf{유전자:} \emph{HSPA1A}, \emph{DNAJB1}, \emph{HSPA9}. 샤페론 \(\REL\)을 15\% 증가.\\
	오차 감소 인자: \(f_3 \approx 0.80\).\\
	\textit{예측:} 응집(HTT‑Q72)이 WT의 \(\le 20\%\)로 감소.
	
	\subsection{층 4 – ECM–인테그린 공명}
	\textbf{유전자:} \emph{ITGAV}, \emph{ITGB3}, \emph{FN1}. 피브로넥틴과의 \(C_{AB}\)를 \(\to 0.97\).\\
	인자: \(f_4 \approx 0.88\).\\
	\textit{예측:} 이동 속도가 WT의 약 125\%로.
	
	\subsection{복합 효과와 주의점}
	\textit{네 층이 독립적으로 실패한다고 가정하면,} 전체 치명적 오차 확률은 개별 오차 감소의 곱이다:
	\[
	\frac{\varepsilon_{\mathrm{LL}}}{\varepsilon_{\mathrm{WT}}} \approx \prod_{i=1}^{4} f_i \approx 0.50.
	\]
	복제 수명이 \(1/\varepsilon\)에 비례한다는 가설 아래, 이는 예를 들어 \textit{in vitro}에서 약 120년에서 약 240년으로의 2배 증가를 의미한다. 이것은 \textbf{상한선}이며; 현실에서는 교차 대화와 자원 공유가 시너지 효과를 감소시키므로 실제 연장은 더 작을 가능성이 높다. 이 예측은 장기 세포 배양 실험에서 명시적으로 테스트 가능하다.
	
	\subsection{제어 조정 회귀 (CCR)}
	인과 관계를 검증하기 위해, 장수 세포는 siRNA/CRISPRi로 일시적으로 되돌려진다(표~\ref{tab:CCR}). 바이오마커는 노화 표현형 쪽으로 이동해야 하며, 세척 또는 siRNA 희석 후 회복되어야 한다.
	
	\begin{table}[H]
		\centering
		\caption{표준 CCR 중재}
		\label{tab:CCR}
		\begin{tabular}{@{}llc@{}}
			\toprule
			\textbf{층} & \textbf{방법} & \textbf{표적 \(\REL\)} \\
			\midrule
			1 & siRNA BRCA1\_LL (50\% KD) & 0.160 \\
			2 & CRISPR 염기 편집 NDUFS1 (I182F) & 0.155 \\
			3 & Tet‑CRISPRi HSPA1A\_LL & 0.170 \\
			4 & siRNA ITGB3\_LL & 0.165 \\
			\bottomrule
		\end{tabular}
	\end{table}
	
	% ============================================================
	% 8. 고급 생물 모듈
	% ============================================================
	\section{고급 생물 모듈}
	\label{sec:modules}
	다음 모듈은 BCLM을 실용적인 생물학적 공학으로 확장하며, 통일된 \(\REL\) 프레임워크로 표현된다.
	
	\subsection{시간 조정}
	단백질 반감기 \(t_{1/2}\)는 \(\REL_{\mathrm{temp}}^{\beta}\)로 스케일한다. 표~\ref{tab:temporal}에 예를 나타낸다.
	
	\begin{longtable}{|c|c|c|}
		\caption{시간 조정 \(\REL\)}\label{tab:temporal}\\
		\hline
		\textbf{단백질} & \(t_{1/2}\) (h) & \(\REL_{\mathrm{temp}}\) \\
		\hline
		\endfirsthead
		\hline
		\textbf{단백질} & \(t_{1/2}\) (h) & \(\REL_{\mathrm{temp}}\) \\
		\hline
		\endhead
		\hline
		p53 & 20 & 0.48 \\
		\(\beta\)-액틴 & 48 & 0.72 \\
		사이클린 B & 1.5 & 0.12 \\
		오르니틴 탈탄산효소 & 0.5 & 0.06 \\
		\hline
	\end{longtable}
	
	\subsection{면역 적합성 행렬}
	MHC–펩타이드 결합은 대수적 적합성 규칙에 의해 지배된다. 표~\ref{tab:mhc}에 예를 나타낸다.
	
	\begin{longtable}{|c|c|c|}
		\caption{MHC–펩타이드 적합성 \(C_{AB}\)}\label{tab:mhc}\\
		\hline
		\textbf{MHC 대립유전자} & \textbf{펩타이드} & \(C_{AB}\) \\
		\hline
		\endfirsthead
		\hline
		\textbf{MHC 대립유전자} & \textbf{펩타이드} & \(C_{AB}\) \\
		\hline
		\endhead
		\hline
		HLA-A*02:01 & FLU NP 265–273 & 0.88 \\
		HLA-A*02:01 & HIV Gag 77–85 & 0.72 \\
		HLA-B*07:02 & EBV LMP2 329–337 & 0.91 \\
		\hline
	\end{longtable}
	
	\subsection{조정 독물 디렉토리}
	특정 억제제는 \(\REL\)을 급격히 감소시킨다. 표~\ref{tab:poisons}에 예를 나타낸다.
	
	\begin{longtable}{|c|c|c|}
		\caption{조정 독물}\label{tab:poisons}\\
		\hline
		\textbf{억제제} & \textbf{표적} & \(\Delta \REL\) \\
		\hline
		\endfirsthead
		\hline
		\textbf{억제제} & \textbf{표적} & \(\Delta \REL\) \\
		\hline
		\endhead
		\hline
		Hg\(^{2+}\) & 티오레독신 환원효소 & \(-80\%\) \\
		뉴클레오시드 유사체 & 바이러스 RT & \(-70\%\) \\
		유기인계 & 아세틸콜린에스테라제 & \(-95\%\) \\
		\hline
	\end{longtable}
	
	\subsection{다중 생물체 조정}
	숙주–마이크로바이옴 상호작용은 \(C_{AB}\)에 의해 정량화된다. 표~\ref{tab:microbiome}에 예를 나타낸다.
	
	\begin{longtable}{|c|c|c|}
		\caption{숙주–마이크로바이옴 적합성}\label{tab:microbiome}\\
		\hline
		\textbf{숙주 수용체} & \textbf{미생물 리간드} & \(C_{AB}\) \\
		\hline
		\endfirsthead
		\hline
		\textbf{숙주 수용체} & \textbf{미생물 리간드} & \(C_{AB}\) \\
		\hline
		\endhead
		\hline
		TLR4 (인간) & LPS (E.\ coli) & 0.91 \\
		AhR (인간) & Indole‑3‑propionate & 0.85 \\
		Ffar2 (생쥐) & 부티르산 & 0.78 \\
		\hline
	\end{longtable}
	
	\subsection{샤페론과 번역 후 변형}
	표~\ref{tab:chaperones}와 \ref{tab:ptm}은 조정 매개변수를 제공한다.
	
	\begin{longtable}{|c|c|c|}
		\caption{샤페론 \(\REL\)}\label{tab:chaperones}\\
		\hline
		\textbf{샤페론} & \(\REL\) & \textbf{기질} \\
		\hline
		\endfirsthead
		\hline
		\textbf{샤페론} & \(\REL\) & \textbf{기질} \\
		\hline
		\endhead
		\hline
		DnaK (Hsp70) & 0.45 & 노출된 소수성 패치 \\
		GroEL/ES (Hsp60) & 0.52 & 폴딩 중간체 \\
		트리거 인자 & 0.38 & 새싱 사슬 \\
		\hline
	\end{longtable}
	
	\begin{longtable}{|c|c|c|}
		\caption{PTM 사이트 \(\Delta \REL\)}\label{tab:ptm}\\
		\hline
		\textbf{변형} & \textbf{콘센서스} & \(\Delta \REL\) \\
		\hline
		\endfirsthead
		\hline
		\textbf{변형} & \textbf{콘센서스} & \(\Delta \REL\) \\
		\hline
		\endhead
		\hline
		N-연결 글리코실화 & Asn‑Xxx‑Ser/Thr & +0.015 \\
		인산화 (PKA) & Arg‑Arg‑Xxx‑Ser & +0.020 \\
		유비퀴틴화 (K48) & 라이신 & –0.030 \\
		\hline
	\end{longtable}
	
	\subsection{유전자 요소 매개변수}
	벡터, 프로모터, 대사산물에 \(\REL\)을 할당한다(표~\ref{tab:elements}).
	
	\begin{longtable}{|c|c|}
		\caption{유전자 요소 \(\REL\)}\label{tab:elements}\\
		\hline
		\textbf{요소} & \(\REL\) \\
		\hline
		\endfirsthead
		\hline
		\textbf{요소} & \(\REL\) \\
		\hline
		\endhead
		\hline
		T7 프로모터 & 0.25 \\
		lac 오페론 프로모터 & 0.08 \\
		ColE1 오리진 (pUC) & 0.12 \\
		rrnB 종결자 & 0.15 \\
		클로람페니콜 내성 & 0.09 \\
		카나마이신 내성 & 0.11 \\
		ATP & 0.32 \\
		NADH & 0.28 \\
		아세틸-CoA & 0.21 \\
		\hline
	\end{longtable}
	
	% ============================================================
	% 9. 탄소 편향 해소와 생체 질량의 반정수 양자화 (새로운 절)
	% ============================================================
	\section{탄소 편향 해소와 생체 질량의 반정수 양자화}
	\label{sec:quantization}
	
	부록 J와 Ferra1.2에서 유도된 보편 질량 사다리는 모든 안정된 물리적 객체에 반정수 양자수 \(N_f\)를 할당하며, 그 질량은
	\begin{equation}
		m = m_e \; q^{N_f}, \qquad q = \left(\frac{3}{2}\right)^{1/3} \approx 1.1447,
		\label{eq:mass_ladder}
	\end{equation}
	을 따른다. 여기서 \(m_e = 0.5109989\)~MeV는 전자 질량이다. 이 사다리는 원래 기본 페르미온과 가벼운 핵에 대해 확립되었으나, 그 유효성은 추가 자유 매개변수 없이 생물학적 거대 분자 클러스터 나아가 전체 세포로까지 확장된다. 이 순수 대수적 규칙은 생물학에서 등방성 중간 매개변수 \(K_{\mathrm{eff}}\)의 필요성을 제거하며, 부록 C2(V.2.2)에서 무기 물질에 대해 확인된 탄소 편향 문제를 완전히 해소한다.
	
	\subsection{탄소–DNA 공명}
	탄소-12는 양자화된 핵 깊이 \(L_{\mathrm{quant}} = 102.0\)(부록 C2)에 위치한다. DNA 코돈(3개 뉴클레오타이드)의 평균 질량은 약 990~Da이며, 이는 반정수 양자수 \(N_f = 129.5\)에 해당한다. 그 차이
	\[
	\Delta = N_f - L_{\mathrm{quant}} = 129.5 - 102.0 = 27.5
	\]
	는 정확히 반정수이다. 따라서 유전 암호는 d‑YPSDC 프로토콜의 탄소 진공 버스와 완벽한 조정 공명 상태에 있다. 정보를 운반하는 중합체는 그 골격을 형성하는 원자와 동일한 이산 질량 격자 위에 직접 고정된다. 경험적 조정은 전혀 필요하지 않으며, 이 정렬은 \(^{12}\mathrm{C}\) 핵과 평균 뉴클레오타이드의 독립적으로 측정된 질량만으로부터 도출된다.
	
	\subsection{리보솜 노드}
	진핵생물 80S 리보솜의 질량은 약 4.3~MDa이다. 이 질량을 사다리 (\ref{eq:mass_ladder})에 투영하면 \(N_f = 191.5\)가 얻어지며, 실제 질량의 이론적 YUCT 이상값으로부터의 편차는 불과 \(-0.91\%\)에 불과하다. 번역의 중심 장치인 리보솜은 이렇게 엄격히 반정수 단계에 위치하여, 해독 과정의 최대 충실도를 보장한다.
	
	\subsection{세포 매크로 클러스터}
	전형적인 인간 세포의 질량은 약 3~ng(건조 중량)이다. 사다리 위에서의 그 위치는 \(N_f = 313.0\)이며, 편차는 \(1.82\%\)이다. 이는 기능적 세포를 생물학적 조정 네트워크의 상위 안정성 한계에 위치시킨다. 다음 반정수 노드로의 상응하는 이동 없이 질량이 더 증가하면, 진공장의 섀넌 엔트로피가 상승하고 조정 기능의 붕괴가 촉발된다.
	
	\subsection{BCLM에 대한 시사점}
	반정수 양자화 규칙은 임의의 생물학적 구조의 절대 질량 척도를 전자 질량과 보편비 \(3/2\)만으로 완전히 결정한다. 이는 질량 의존적 거시적 특성(대사 스케일링, 오차율 등)을 예측할 때 개별 매개변수 \(\REL\)을 사용할 필요성을 제거한다. 대신 양자수 \(N_f\)를 적합성 행렬에 직접 사용할 수 있다.
	\[
	\Delta_{AB} = |N_f(A) - N_f(B)|,
	\]
	동일한 가우스 창 \(C_{AB} = \exp[-(\Delta_{AB} - n_{AB})^2/2\sigma^2]\) 및 \(\sigma=0.20\)과 함께. 이로써 생물학적 라그랑지안은 물리적 진공과의 매개변수 없는 연결을 획득하여, 클라이버의 법칙과 페르미온 질량 계층을 통일한다.
	
	\subsection{예측}
	\begin{enumerate}
		\item 모든 기능적 생체 고분자 복합체(리보솜, 프로테아솜, 스플라이소솜)는 정확한 반정수 \(N_f\)의 \(\pm 0.25\) 반정수 단위 내에 질량을 가져야 한다. 알려진 복합체에 대한 체계적 조사가 진행 중이며, 예비 결과는 \(0.25\)를 초과하는 편차가 충실도나 안정성 저하와 상관관계가 있음을 보여준다.
		\item 살아있는 세포의 최대 크기는 노드 \(N_f \approx 313.5\)에 의해 제한되며, 이 질량을 초과하는 세포는 분열하거나 노화 상태로 들어가야만 한다. 이는 진핵생물 전반에서 관찰되는 세포 크기 상한에 대한 기작적 설명을 제공한다.
		\item \(\Delta = 27.5\)에서의 탄소–DNA 공명은 다른 골격 화학을 사용하는 대체 유전 시스템이 동등한 복제 충실도를 달성하려면 동일한 정밀도로 \(3/2\) 사다리와 일치해야 함을 시사한다. 이는 합성 이종 핵산(XNA)으로 검증할 수 있다.
	\end{enumerate}
	
	이렇게 반정수 사다리는 부록 C의 비생물적 조정 효율과 섹션 \ref{sec:codons}–\ref{sec:enzymes}에 표로 제시된 생물학적 효율 사이의 결여된 연결 고리를 제공한다. 이는 보편 지수 \(\beta = 2/3\)이 단순한 수비학적 호기심이 아니라, 무생물과 생물 모두를 조직하는 이산적 층상 진공의 구조적 상수임을 확증한다.
	
	% ============================================================
	% 10. 검증 현황 및 제안된 실험
	% ============================================================
	\section{검증 현황 및 제안된 실험}
	\label{sec:validation}
	
	\subsection{기존의 독립적 확인}
	보편 지수 \(\beta = 2/3\)는 여러 비생물학적 시스템(결합 에너지, 상전이, 우주 마이크로파 배경)에서 확고하게 확립되었다. 생물학 내에서, 다음 테스트는 순환 논리 없이 지지를 제공한다:
	\begin{itemize}
		\item \textbf{대사 스케일링:} 지수 \(b = \beta d_f/2\)는 관측된 \(0.67\)(단순 생물)과 \(0.74\)(포유류)를 단일 \(\beta\)로 설명하며, 1016종 식물과 379종 포유류의 분석에 의해 확인되었다(부록 D).
		\item \textbf{신경 동기화:} 훈련 유도에 의한 \(\REL\)의 1.81배 증가는 예측된 오차 감소 \(\varepsilon \propto \REL^{-0.67}\)와 일치한다(부록 D).
	\end{itemize}
	돌연변이율 데이터는 법칙의 검증이 아니라(순환 논리가 될 것임) \(\alpha_{\mathrm{repair}}\)의 보정 역할을 한다. DNA 복구 능력의 독립적 측정이 그 관계를 완전히 검증하는 데 필요하다.
	
	\subsection{제안된 중요 실험}
	\begin{enumerate}
		\item \textbf{단백질-단백질 복합체의 \(\REL\) 직접 측정:} 알려진 \(\REL\)을 가진 두 단백질을 발현시키고, \(K_d\)를 측정하며, 식(\ref{eq:Kd})의 예측과 비교한다. 이는 적합성 규칙을 직접 테스트할 것이다.
		\item \textbf{번역 오차율 대 코돈 \(\REL\):} 두 개의 GFP 변이체(최적화 대 무작위 코돈)를 합성하고, 질량 분석을 통해 형광과 오삽입율을 비교한다.
		\item \textbf{온도가 돌연변이율에 미치는 영향:} 동일한 복구 유전자 발현 배경 하에서, 효모의 \(25^\circ\text{C}\)와 \(37^\circ\text{C}\)에서 \(\mu\)를 평가한다.
		\item \textbf{장수 파일럿 분석:} 네 층 구성체를 인간 iPS 유래 심근세포에 도입하고, 1년 동안 HPRT 돌연변이, MitoSOX, 봉입체, 누적 배가수를 모니터링한다.
	\end{enumerate}
	
	% ============================================================
	% 11. 결론
	% ============================================================
	\section{결론}
	이 버전의 BCLM은 YUCT를 생물학에 적용하기 위한 자족적이고 엄격하며 투명한 프레임워크를 제공한다. 핵심량 \(\REL\)은 작동적으로 정의되고 공개 데이터로부터 균일하게 계산된다. 적합성 규칙은 대규모 실험 데이터베이스에 대해 보정되었으며, 그 자유 매개변수는 고정되었다. 예측은 명확하게 기술되고, 가정과 한계는 명시적으로 인정된다. 따라서 BCLM은 YUCT 라그랑지안의 검증 가능한 생물학적 섹터를 형성한다.
	
	\vspace{0.5cm}
	\noindent\textbf{보충 자료:} 보정 데이터, Python 스크립트, 전체 SKEMPI 표는 \url{https://github.com/Alexey-Yakushev-YUCT/YPSDC}에서 이용 가능하다.
	
	\begin{thebibliography}{99}
		\bibitem{codon_usage_db} Codon Usage Database, \url{https://www.kazusa.or.jp/codon/}.
		\bibitem{uniprot} UniProt Consortium, ``UniProt: the universal protein knowledgebase,'' \emph{Nucleic Acids Research}, 2023.
		\bibitem{brenda} BRENDA Enzyme Database, \url{https://www.brenda-enzymes.org/}.
		\bibitem{pdb} H.M. Berman et al., ``The Protein Data Bank,'' \emph{Nucleic Acids Res.} 28, 235–242 (2000).
		\bibitem{skempi} J. Jankauskaitė et al., ``SKEMPI 2.0,'' \emph{Nucleic Acids Res.} 47, D535–D541 (2019).
		\bibitem{bergeron2023} L.\,A.\,Bergeron \emph{et al.}, ``Evolution of the germline mutation rate across vertebrates,'' \emph{Nature}, 615, 285–291 (2023).
		\bibitem{sharp1987} P.\,M.\,Sharp and W.\,H.\,Li, ``The codon adaptation index,'' \emph{Nucleic Acids Res.}, 15, 1281–1295 (1987).
		\bibitem{kastritis2011} P.\,L.\,Kastritis \emph{et al.}, ``On the binding affinity of macromolecular interactions,'' \emph{Curr. Opin. Struct. Biol.}, 21, 50–61 (2011).
		\bibitem{bareven2011} A.\,Bar‑Even \emph{et al.}, ``The moderately efficient enzyme,'' \emph{Biochemistry}, 50, 4402–4410 (2011).
		\bibitem{AppendixA} A.\,V.\,Yakushev, ``Appendix A: Practical Guide to Using YUCT V35.0,'' in \emph{YUCT}, Zenodo (2026).
		\bibitem{AppendixD} A.\,V.\,Yakushev, ``Appendix D: Biological Predictions,'' in \emph{YUCT}, Zenodo (2026).
		\bibitem{AppendixP} A.\,V.\,Yakushev, ``Appendix P: Temperature Dependence of Gravity,'' in \emph{YUCT}, Zenodo (2026).
	\end{thebibliography}
	
\end{document}